\scanpage{97}
% archive sheet number: 96
\begin{center}
Filtration et poids des espaces de cohomologie des variétés algébriques

par A.GROTHENDIECK
\end{center}

0. \uline{Introduction}. On a remarqué, depuis \deleted{longtemps}\inserted{l'invention de la cohomologie $\handwritten{\ell}$-adique \handwritten{\unsure{([\quad]|[\quad])}}}, que\deleted{\unsure{,|x}} pour une variété algébrique $X$ définie sur un corps $k$, de clôture algébrique $\bar{k}$, les opérations naturelles du groupe de Galois $\deleted{\unsure{\operatorname{Gal}k|G_k|G(k)}}\handwritten{\pi}=\operatorname{Gal}(\bar{k}/k)$ sur la cohomologie $\handwritten{\ell}$-adique $H(X_{\bar{k}},\underline{Z}_{\handwritten{\ell}})$ de la variété ``géométrique'' $X_{\bar{k}}$ définissait sur cette \deleted{dernière}\inserted{\handwritten{cohomologie}} une structure supplémentaire remarquable, donnant lieu notamment aux remarquables conjectures de Tate \handwritten{[\quad]}. Il semble probable que ces dernières soient n\deleted{\unsure{o|e}}\handwritten{e}t\deleted{\unsure{ablem|able}}\handwritten{tem}ent plus éloignées d\handwritten{e}\deleted{leur solution} leur solution que les conjectures de Weil \handwritten{[\quad,\quad]}, qui, comme on sait, peuvent s'interpréter également en termes d'actions galoisiennes sur des groupes de cohomologie $\handwritten{\ell}$-adiques. Notre propos dans le présent travail est de montrer comment on peut, des conjectures de Weil et (le cas échéant) de la résolution des singularités, tirer des conséquences de nature ``géométrique'' en apparence, i.e. relatives à un corps de base algébriquement clos. Le principe de cette méthode\inserted{\deleted{(qui remonte à 1964)}}, qui consiste donc essentiellement à ramener un énoncé géométrique au cas d'un corps de base fini, est expliqué au N°1, et des exemples d'application de cette méthode sont donnés dans les numéros 2 à 7. \deleted{\unsure{Le pr}} L'idée de ces raisonnements, liés à\deleted{\unsure{ux|la}}\inserted{la} considération\deleted{s}\inserted{systématique} des ``poids'' des modules galoisiens provenant de la cohomologie $\handwritten{\ell}$-adique, remonte à 1964; des exemples \deleted{du même principe}\inserted{\handwritten{analogues}} se trouvent développés dans SGA 7\handwritten{(}I et II\handwritten{)}; dans le séminaire oral\deleted{e} était\inserted{\handwritten{développé}} également l'\unsure{exempbe|exemple} du n°2 du présent travail.

Bien entendu, la portée des ``résultats'' de notre travail est nécessairement \deleted{limitée}\inserted{\handwritten{sujette à caution}}, puisqu'ils dépendent de conjectures non prouvées à l'heure actuelle, savoir les conjectures de Weil et la résolution des singularités (sous la forme forte de Hironaka) en caractéristique
